\section{Does the parallel path compute the same thing?}
\label{sec:correct}

Before timing anything we checked what each rung produces. Table~\ref{tab:correct}
summarises the cases; the raw records are in \texttt{e1\_correctness.jsonl} and
\texttt{e0\_traces.jsonl}.

\begin{table}[t]
\caption{Output of each rung, compared by a hash of \texttt{filelist.txt} and of every
listed \texttt{<inode>.txt}. Same letter = byte-identical output.}
\label{tab:correct}
\small
\begin{tabular}{@{}lcccc@{}}
\toprule
Case & (a) & (b) & (c) & (d)/(e) \\
\midrule
A: aligned \texttt{O\_DIRECT}, 60 files & X & X & X & X \\
B: unaligned buffered requests & X & Y & X & Y \\
C: tmpfs (no FIEMAP), 40 files & 40 & 40 & 40 & \textbf{0} files \\
C: 9p (no FIEMAP), 30 files & 30 & 30 & 30 & \textbf{0} files \\
D: unrelated \texttt{notes.txt} in cwd & kept & -- & kept & \textbf{re-sorted} \\
E: 10{,}000 traced files & -- & -- & 1.1\,s & 25.8\,s \\
E: 10{,}001 traced files & -- & -- & 0.9\,s & \textbf{hangs} \\
\bottomrule
\end{tabular}
\end{table}

\paragraph{The extent lookup is extra work whose result nothing reads.}
Upstream's analysis phase does not read extents (\S\ref{sec:background}). The
parallel path does, writes them to a temporary directory that is never removed,
sorts them there, and drops them: \texttt{merge.py}, \texttt{hotness.sh} and both
migration scripts read only \texttt{filelist.txt} and the per-inode request files,
and the migration scripts look the extents up again themselves. The only effect of
the lookup on the result is a filter: a file whose FIEMAP call fails is left out of
\texttt{filelist.txt}. On a file system without FIEMAP the parallel path therefore
lists no files at all (case~C), while upstream lists all of them. This matters for
the fork's own motivation, which names network file systems as the target; we
observed it on tmpfs and on WSL's 9p file system.

\paragraph{The two code paths disagree on unaligned requests.}
Phase~1 rounds the end of each request up to a block boundary. Upstream computes
\texttt{end += end \% 4096 - 1}, which for a 100-byte read at offset~0 gives
\texttt{end} = 197; the fork computes \texttt{end += 4096 - ((end+1) \% 4096) - 1},
which gives 4094. Neither gives 4095. Both paths of the fork use the second formula,
so (b) and (d) agree with each other and differ from upstream on buffered workloads
with unaligned requests (case~B). Upstream's own evaluation issues only 128\,KiB
\texttt{O\_DIRECT} requests, which never reach either formula.

\paragraph{The parallel engine deadlocks above its queue size.}
\texttt{EnhancedParallel\allowbreak Processor} puts every task into a bounded queue
(\texttt{queue\_size} = 10{,}000) before it starts the worker threads, so the
10{,}001st \texttt{put} blocks forever (case~E). When \texttt{--workers} is omitted,
auto-tuning lowers the bound to 1{,}000 if less than 4\,GB of memory is available,
so the limit can be a thousand traced files.

\paragraph{The parallel path rewrites unrelated files.}
Its third step sorts every \texttt{*.txt} file in the working directory, including
files that are not analysis output (case~D).

\paragraph{Two of the fork's own scripts do not run.} The fork's \texttt{analyze.sh},
the entry point that chains tracing and analysis, fails \texttt{bash -n} with a syntax
error the fork introduced (\texttt{if ! ./trace.sh "\$PROCESS\_NAME" \&; then}, line
76); upstream's version parses. Its \texttt{overhead\_profiling\_benchmark.py} stops
at an import of a class that was later renamed. The two speedup benchmarks, one of
which produced the 12.74$\times$ figure, run.

\paragraph{The trace itself loses and corrupts events.}
These problems are upstream's and affect every rung equally.
\begin{itemize}
\item BCC's \texttt{trace} flushes its standard output only after printing its header,
  and upstream stops the tracers with \texttt{kill} (SIGTERM); the \texttt{kill -INT}
  in \texttt{analyze.sh} goes to \texttt{trace.sh}, which has already exited, and
  background jobs of a non-interactive shell ignore SIGINT anyway. The last
  block-buffered chunk of events is lost: 59 of the 8{,}192 reads of upstream's
  \texttt{read\_seq} benchmark (the end of the file), and all events of a workload
  with a single read. Setting \texttt{PYTHON\allowbreak UNBUFFERED=1} for the tracers recovers
  the tail but slows the consumer enough that BCC drops samples at high event rates
  (it reported 3{,}739 lost samples during \texttt{tar}, about 13\% of the events).
\item \texttt{trace.sh} prints the file position with \texttt{\%u}, a 32-bit format
  for the 64-bit \texttt{loff\_t}: a read at 5\,GiB is recorded at 1\,GiB.
\item Trace lines carry an inode number but no device number. The traced processes
  also read shared libraries and locale files on the root file system (between 2
  and 89 distinct such inodes per trace in our six workloads), and any of those numbers that
  also exists under \texttt{/mnt} is attributed to the wrong file.
\item When the two tracers start at the same time on a kernel without a usable
  header tree (as in WSL2), BCC's fallback to the \texttt{kheaders} module races
  between them and one tracer fails to compile, silently losing all reads or all
  writes. We avoid this by pointing \texttt{BCC\_KERNEL\_SOURCE} at headers
  extracted once.
\end{itemize}

\paragraph{The rewrite.}
Rung (c) produced the same bytes as upstream in every case above, including
case~B, because it reproduces upstream's formula; we kept the bug on purpose so that
the timing comparison is between equal outputs.
